\section{Open Questions}

This version leaves several questions intentionally open.

\begin{enumerate}[label=\textbf{Q\arabic*.}]
    \item \textbf{Second order or third order?}
    The second-order expansion identifies risk correction terms in equilibrium
    equations. However, a fully solved stochastic-volatility DSGE model usually
    requires third-order perturbation for volatility shocks to affect decision
    rules as independent state variables. The next step is to decide whether the
    paper needs a full third-order solution or only an accounting
    representation.

    \item \textbf{Where should wedges be defined?}
    This note defines wedges in the reduced-form IS curve, NKPC, and Taylor
    rule. A deeper BCA exercise could instead define wedges at the level of
    household Euler equations, firm pricing FOCs, and resource constraints.

    \item \textbf{Which pricing model pins down \(\lambda_\pi\)?}
    The sign of the price-setting wedge depends on the nonlinear pricing
    environment. Calvo and Rotemberg pricing can imply different uncertainty
    transmission mechanisms, so the markup wedge should not be interpreted
    before choosing the pricing microfoundation.

    \item \textbf{Should natural-rate risk be absorbed into \(\omega^x_t\)?}
    One can define a risk-adjusted natural rate and keep the IS wedge smaller,
    or define the natural rate at first order and place all risk corrections in
    \(\omega^x_t\). The latter is cleaner for accounting; the former may be
    cleaner for structural interpretation.

    \item \textbf{How will this connect to the empirical uncertainty factors?}
    This note has one volatility state \(q_t\). The larger research project has
    empirical macro and financial uncertainty factors. The future SOE version
    will need a factor-loading structure that maps those empirical factors into
    multiple structural wedge volatilities.
\end{enumerate}
